\documentclass[12pt]{book}
\usepackage{precalc1}
\setcounter{chapter}{2}
\begin{document}

\chapter{Polynomials and Their Shape}

A polynomial of high degree looks intimidating in expanded form, a long sum of
powers with no obvious shape. The global shape of its graph is nonetheless forced by
a small amount of local data: how the function behaves at the far left and far right,
and what happens at each root. The factored form displays this data. Each factor
names a root, the exponent on that factor governs whether the graph crosses the axis
or turns back from it, and the leading term governs both ends.

\section{What a Polynomial Is}

\begin{definition}{Polynomial}
A \emph{polynomial} is a function of the form
\[
f(x)=a_n x^n+a_{n-1}x^{n-1}+\cdots+a_1 x+a_0,
\]
where the coefficients $a_i$ are real numbers and $a_n\neq 0$. The \emph{degree} is
$n$, the highest power present; $a_n$ is the \emph{leading coefficient} and $a_0$ is
the constant term.
\end{definition}

Lines are degree-one polynomials and parabolas are degree-two. The expanded form
above is convenient for evaluation and for reading the constant term $a_0=f(0)$, but
it conceals the shape. The factored form reveals it.

\begin{keyidea}
The expanded form of a polynomial is built for arithmetic; the factored form is
built for shape. Where the graph crosses the axis, where it turns back, and where it
goes at the two ends are all read from the factored form together with the leading
term.
\end{keyidea}

\section{End Behavior: the Two Ends}

Far from the origin, the highest-degree term dominates every other term, because a
high power outgrows all lower powers once $|x|$ is large. So the two ends of a
polynomial's graph behave like the graph of its leading term $a_n x^n$ alone. Two
features of that term decide the ends: the parity of the degree $n$ and the sign of
the leading coefficient $a_n$.

\begin{keyidea}
\textbf{End behavior of} $f(x)=a_n x^n+\cdots$
\begin{itemize}\setlength{\itemsep}{2pt}
  \item \emph{Even degree.} Both ends point the same way: up if $a_n>0$, down if
        $a_n<0$. The graph resembles a parabola at large scale.
  \item \emph{Odd degree.} The ends point opposite ways: down-left to up-right if
        $a_n>0$, up-left to down-right if $a_n<0$. The graph resembles a cubic at
        large scale.
\end{itemize}
\end{keyidea}

\begin{center}
\begin{tikzpicture}[scale=0.62]
  % even, positive
  \begin{scope}
    \draw[->] (-2.4,0)--(2.4,0); \draw[->] (0,-1)--(0,3.6);
    \draw[pcblue,thick,domain=-1.8:1.8,smooth] plot (\x,{0.6*\x*\x*\x*\x-0.3});
    \node at (0,-1.8) {\scriptsize even, $a_n>0$};
  \end{scope}
  % even, negative
  \begin{scope}[xshift=6cm]
    \draw[->] (-2.4,0)--(2.4,0); \draw[->] (0,-3.6)--(0,1);
    \draw[pcred,thick,domain=-1.8:1.8,smooth] plot (\x,{-0.6*\x*\x*\x*\x+0.3});
    \node at (0,-4.4) {\scriptsize even, $a_n<0$};
  \end{scope}
  % odd, positive
  \begin{scope}[xshift=12cm]
    \draw[->] (-2.4,0)--(2.4,0); \draw[->] (0,-2.6)--(0,2.6);
    \draw[pcblue,thick,domain=-1.7:1.7,smooth] plot (\x,{0.5*\x*\x*\x});
    \node at (0,-3.4) {\scriptsize odd, $a_n>0$};
  \end{scope}
  % odd, negative
  \begin{scope}[xshift=18cm]
    \draw[->] (-2.4,0)--(2.4,0); \draw[->] (0,-2.6)--(0,2.6);
    \draw[pcred,thick,domain=-1.7:1.7,smooth] plot (\x,{-0.5*\x*\x*\x});
    \node at (0,-3.4) {\scriptsize odd, $a_n<0$};
  \end{scope}
\end{tikzpicture}
\end{center}

\begin{example}[reading the ends]
The polynomial $f(x)=-2x^5+7x^2-4$ has odd degree $5$ and negative leading
coefficient. The ends therefore run up on the left and down on the right. No
factoring is needed; the leading term $-2x^5$ settles both ends by itself.
\end{example}

\begin{pitfall}
End behavior depends only on the leading term, never on the constant term or the
middle coefficients. A large constant shifts the graph vertically but does not change
where the ends go. Reading $a_0$ for end behavior is a common error; read $a_n x^n$.
\end{pitfall}

\section{Roots and the Factored Form}

A \emph{root} (or zero) of $f$ is an input $r$ with $f(r)=0$, a point where the graph
meets the $x$-axis. The factored form lists the roots directly:
\[
f(x)=a_n (x-r_1)^{m_1}(x-r_2)^{m_2}\cdots (x-r_k)^{m_k}.
\]
Each $r_i$ is a root, and the exponent $m_i$ is its \emph{multiplicity}, the number
of times that factor is repeated. The multiplicities sum to the degree $n$.

\begin{definition}{Multiplicity}
The \emph{multiplicity} of a root $r$ is the exponent on the factor $(x-r)$ in the
factored form. A root with multiplicity one is \emph{simple}; higher multiplicities
are \emph{repeated}.
\end{definition}

Multiplicity controls how the graph meets the axis at the root, and the rule is a
parity rule, the same parity idea that governed the ends.

\begin{keyidea}
\textbf{Behavior at a root of multiplicity} $m$:
\begin{itemize}\setlength{\itemsep}{2pt}
  \item \emph{$m$ odd.} The graph \emph{crosses} the axis; the function changes
        sign there. ($m=1$: a straight crossing. $m=3,5,\dots$: a crossing that
        flattens against the axis.)
  \item \emph{$m$ even.} The graph \emph{touches} the axis and turns back; the
        function does \emph{not} change sign there. ($m=2$: a parabola-like bounce.)
\end{itemize}
\end{keyidea}

\begin{visual}
At a root, the factor $(x-r)^m$ changes sign exactly when $m$ is odd, because a
negative number raised to an odd power stays negative while an even power makes it
positive. Crossing the axis is a sign change; touching and turning back is not. The
multiplicity's parity is the whole story.
\end{visual}

\begin{center}
\begin{tikzpicture}[scale=0.7]
  % simple root: crosses
  \begin{scope}
    \draw[->] (-2,0)--(2,0); \draw[->] (0,-1.8)--(0,1.8);
    \draw[pcblue,thick,domain=-1.5:1.5,smooth] plot (\x,{\x});
    \fill[pcred] (0,0) circle (2pt);
    \node at (0,-2.6) {\scriptsize $m=1$: crosses};
  \end{scope}
  % double root: touches
  \begin{scope}[xshift=5.5cm]
    \draw[->] (-2,0)--(2,0); \draw[->] (0,-0.6)--(0,2.6);
    \draw[pcgreen,thick,domain=-1.45:1.45,smooth] plot (\x,{\x*\x});
    \fill[pcred] (0,0) circle (2pt);
    \node at (0,-1.4) {\scriptsize $m=2$: touches};
  \end{scope}
  % triple root: flat crossing
  \begin{scope}[xshift=11cm]
    \draw[->] (-2,0)--(2,0); \draw[->] (0,-1.8)--(0,1.8);
    \draw[pcblue,thick,domain=-1.2:1.2,smooth,samples=60] plot (\x,{\x*\x*\x});
    \fill[pcred] (0,0) circle (2pt);
    \node at (0,-2.6) {\scriptsize $m=3$: flat crossing};
  \end{scope}
\end{tikzpicture}
\end{center}

\begin{example}[reading roots and multiplicities]
For $f(x)=(x+2)(x-1)^2(x-3)$, the roots are $-2$, $1$, and $3$. The roots $-2$ and
$3$ are simple (odd multiplicity), so the graph crosses there. The root $1$ has
multiplicity two (even), so the graph touches the axis at $x=1$ and turns back
without crossing. The degree is $1+2+1=4$, even, with positive leading coefficient,
so both ends point up.
\end{example}

\demo{Build a polynomial by placing roots and setting each multiplicity, in the form
$y=(x-r_1)^{m_1}\cdots(x-r_n)^{m_n}$, then fit a hidden target. The fit score weighs
root locations, multiplicities, and overall agreement. Watch how an even
multiplicity makes the curve bounce off the axis while an odd multiplicity sends it
through.}{https://bobovski66.github.io/Precalculus/demos/polynomial\_roots.html}

\section{The Sign Chart}

The roots are the only places a polynomial can change sign, because a polynomial is
continuous and changes sign only by passing through zero. Between consecutive roots
the function holds a single sign throughout. This is the observation that turns local
data into a global sketch: mark the roots on a number line, and the line breaks into
intervals on each of which $f$ is entirely positive or entirely negative. One test
point per interval reveals which.

\begin{keyidea}
\textbf{Sign chart procedure.}
\begin{enumerate}\setlength{\itemsep}{2pt}
  \item Factor $f$ and mark all roots on a number line; they divide it into intervals.
  \item Pick one test point in each interval and evaluate the sign of $f$ there
        (the sign of the product of factors).
  \item Label each interval $+$ or $-$. The sign can change only at a root, and only
        at a root of odd multiplicity.
\end{enumerate}
\end{keyidea}

\begin{example}[building a sign chart]
Analyze $f(x)=(x+2)(x-1)^2(x-3)$.

The roots $-2,1,3$ divide the line into four intervals. Evaluate the sign of each
factor at a test point in each interval; the factor $(x-1)^2$ is never negative, so
it never affects the sign.
\begin{center}
\renewcommand{\arraystretch}{1.3}
\begin{tabular}{@{}lccccc@{}}
\toprule
Interval & test $x$ & $(x+2)$ & $(x-1)^2$ & $(x-3)$ & sign of $f$\\
\midrule
$(-\infty,-2)$ & $-3$ & $-$ & $+$ & $-$ & $+$\\
$(-2,1)$       & $0$  & $+$ & $+$ & $-$ & $-$\\
$(1,3)$        & $2$  & $+$ & $+$ & $-$ & $-$\\
$(3,\infty)$   & $4$  & $+$ & $+$ & $+$ & $+$\\
\bottomrule
\end{tabular}
\end{center}
The sign stays negative across $x=1$, confirming the even multiplicity: no sign
change there. It flips at $x=-2$ and at $x=3$, the simple roots.
\end{example}

\begin{center}
\begin{tikzpicture}[scale=1.0]
  \draw[thick] (-4.5,0)--(4.5,0);
  \foreach \x/\lab in {-3/-2,0/1,3/3}{
    \fill[pcred] (\x,0) circle (2.4pt);
  }
  \node[below] at (-3,-0.15) {\scriptsize $-2$};
  \node[below] at (0,-0.15) {\scriptsize $1$};
  \node[below] at (3,-0.15) {\scriptsize $3$};
  \node[pcgreen] at (-4,0.35) {\scriptsize $+$};
  \node[pcred] at (-1.5,0.35) {\scriptsize $-$};
  \node[pcred] at (1.5,0.35) {\scriptsize $-$};
  \node[pcgreen] at (4,0.35) {\scriptsize $+$};
\end{tikzpicture}
\end{center}

\begin{visual}
The sign chart is a number line with the roots marked as special points and a sign
assigned to each interval between them. It records, at a glance, where the function
lives above the axis and where below. The graph must obey it: positive intervals
hold the curve above the axis, negative intervals below, and the roots pin it to the
axis.
\end{visual}

\section{Assembling the Global Sketch}

End behavior fixes the two ends. The sign chart fixes which side of the axis the
curve occupies on every interval. Multiplicity fixes how it meets each root. These
three pieces of local data determine the global shape up to the heights of the turns,
which precalculus does not pin down exactly but calculus does. A faithful sketch
needs only the three pieces.

\begin{keyidea}
\textbf{Sketching a polynomial from its factored form.}
\begin{enumerate}\setlength{\itemsep}{2pt}
  \item Read the degree and leading coefficient; draw the two ends.
  \item Mark the roots; at each, cross (odd multiplicity) or touch (even).
  \item Use the sign chart to keep the curve on the correct side of the axis between
        roots.
  \item Connect smoothly, consistent with the ends and the sign chart.
\end{enumerate}
\end{keyidea}

\begin{example}[a full sketch]
Sketch $f(x)=(x+2)(x-1)^2(x-3)$.

\emph{Ends.} Degree $4$, positive leading coefficient: both ends up.

\emph{Roots.} Cross at $-2$, touch at $1$, cross at $3$.

\emph{Sign chart.} Positive on $(-\infty,-2)$, negative on $(-2,3)$ except pinned to
zero at $x=1$, positive on $(3,\infty)$.

The curve descends from the upper left, crosses down at $-2$, runs below the axis,
rises to touch zero at $x=1$ and falls back below, then crosses up at $3$ and climbs
to the upper right. The touch at $x=1$ is the visual signature of the double root.
\end{example}

\begin{center}
\begin{tikzpicture}[scale=0.95]
  \draw[->] (-3.6,0)--(4.4,0) node[right] {$x$};
  \draw[->] (0,-3.4)--(0,3.6) node[above] {$y$};
  \draw[pcblue,thick,domain=-2.35:3.4,smooth,samples=160]
    plot (\x,{0.1*(\x+2)*(\x-1)*(\x-1)*(\x-3)});
  \foreach \x in {-2,1,3}\fill[pcred] (\x,0) circle (2.2pt);
  \node[below left] at (-2,0) {\scriptsize $-2$};
  \node[above right] at (1,0.05) {\scriptsize $1$};
  \node[below right] at (3,0) {\scriptsize $3$};
\end{tikzpicture}
\end{center}

\begin{pitfall}
The sign chart constrains which side of the axis the curve is on, not how high it
rises. Without calculus the exact heights of the peaks and valleys are not
determined, so sketches should show the correct \emph{shape} (crossings, touches,
sides, ends) and not claim exact turning-point heights.
\end{pitfall}

\section{Exercises}

\begin{exercises}
\item State the end behavior of each polynomial from its leading term alone:
(a) $f(x)=3x^4-x+5$; (b) $f(x)=-x^3+2x^2$; (c) $f(x)=-4x^6+x^5$;
(d) $f(x)=x^5-3x^2+1$.

\item Explain why the constant term and middle coefficients have no effect on end
behavior, referring to how powers compare when $|x|$ is large.

\item For $f(x)=(x+1)(x-2)^3(x-4)$, list the roots with their multiplicities, state
the degree, and say at each root whether the graph crosses or touches.

\item Using the parity of the exponent, explain why a factor $(x-r)^m$ changes sign
at $r$ when $m$ is odd and does not when $m$ is even.

\item Build a complete sign chart for $f(x)=(x+3)(x-1)(x-2)$, with a test point in
each interval and the sign of every factor recorded.

\item Build a sign chart for $f(x)=-(x+1)^2(x-3)$. Identify where the sign does and
does not change, and connect each to a multiplicity.

\item Sketch $f(x)=(x-1)(x+2)^2$ from its factored form: state the ends, mark the
roots with crossing or touching behavior, use the sign chart, and draw a shape
consistent with all three.

\item Sketch a degree-five polynomial with a simple root at $-2$, a triple root at
$0$, a simple root at $3$, and positive leading coefficient. Describe the behavior at
the triple root in words.

\item A polynomial has even degree and negative leading coefficient, with simple
roots at $-1$ and $4$ and a double root at $2$. Determine the lowest possible degree,
then sketch a graph consistent with all stated features.

\item A polynomial is positive on $(-\infty,-1)$, negative on $(-1,2)$, and positive
on $(2,\infty)$, with no other sign changes. Determine the multiplicity parity at
each of $-1$ and $2$, and give a possible factored form of lowest degree.

\item Two students draw graphs of $f(x)=(x-1)^2(x-4)$ that differ only in the height
of the turning point between the roots. Explain why both can be correct sketches and
what additional tool would be needed to fix the height.

\item Open the Roots and Multiplicity Lab. Build a polynomial of degree at least four
with at least one repeated root, record its factored form and the fit score against a
target, and write one sentence on how changing a multiplicity from even to odd
altered the graph at that root.
\end{exercises}

\end{document}
